% Abstract — JNE limit ~250 words.  Structured paragraph (no labels).
Peripheral nerve stimulation is a rapidly expanding therapeutic
modality, but optimising electrode geometry and stimulus waveforms
for fascicle-specific selectivity remains computationally expensive:
state-of-the-art pipelines wrap the NEURON simulator from Python and
require thousands of forward simulations per design candidate. We
present \emph{jaxon}, a fully differentiable reimplementation of the
canonical peripheral nerve fiber models---McIntyre--Richardson--Grill
(MRG) myelinated axons, and the unmyelinated Sundt, Sweeney and Rattay
C-fibers---in JAX. Built on top of the \texttt{jaxley} differentiable
Hodgkin--Huxley framework, jaxon adds a batched coupled-cable
backward-Euler integrator that simulates hundreds of fibers in
parallel on a single GPU and exposes end-to-end gradients with respect
to per-contact stimulus amplitudes and arbitrary waveforms. We
validate jaxon against NEURON on strength--duration, conduction
velocity, and intracellular-trace agreement across all four fiber
models, reproducing established benchmark behaviour to within
\needdata{X}\,\% of NEURON. We then benchmark throughput against
\texttt{pyfibers} and find a \needdata{Y}-fold reduction in wall time
per forward simulation at $N=200$ fibers. As a demonstration of the
differentiability advantage, we use jaxon to optimise spatially
selective fascicular stimulation across $N=100$ randomised Hussain-style
multi-fascicle nerve realisations, achieving a mean selectivity index
of \needdata{0.97 placeholder---update from Phase 3 sweep results}
with rectangular pulses and showing that warm-started waveform
optimisation refines---without destabilising---the rectangular optimum.
jaxon is open-source and designed to slot into existing peripheral nerve
stimulation pipelines as a drop-in replacement for the NEURON forward
solver, with the additional capability of gradient-based optimisation.
